% C4 scored at R=5 by half-filling. Tie. Tracking collapses.
\documentclass[11pt]{article}
\usepackage[margin=1.05in]{geometry}
\usepackage{amsmath,amssymb,amsthm}
\usepackage[colorlinks=true,linkcolor=blue,citecolor=blue,urlcolor=blue]{hyperref}

\newtheorem*{admission}{Admission}

\title{\textbf{C4 Scored at Radius $5$:\\
the Parabola and a Linear Weight Tie}}
\author{Robert W.\ McGwier\\
\small CTO, Cohere Technology Group\\
\small GitHub: \texttt{n4hy}}
\date{15 August 2026 --- Version 1}

\begin{document}
\maketitle

\begin{abstract}
\noindent
Paper~30 could not score C4: $E<0$ filling flipped on $25$ of
$30$ hops. This note occupies the lowest $V/2$ eigenstates, so
$n_{\mathrm{occ}}$ cannot flip.

$N=14$, $R=5$, thirty locked surface hops, $\varepsilon=0.01$,
$n_{\mathrm{flip}}=0$.
$R_{\mathrm{CHM}}=0.99951$, $R_{\mathrm{lin}}=0.99945$,
gap $6\times 10^{-5}$. C4 is scored: \textbf{TIE}.
$\lvert\rho\rvert\sim 0.03$: the first law does not track.

Auditor $N=16$, $R=5$, fifteen hops: $0.889$ versus $0.892$,
again a dead heat; flat loses ($0.997$).

The parabola is not selected at $R=5$. Half-filling buys a
score and loses the $E<0$ tracking of Papers~26--29.
Not Planck. Not Einstein.
\end{abstract}

\section{Numbers}

\begin{center}
\begin{tabular}{lcccc}
\hline
 & $n$ & $R_{\mathrm{CHM}}$ & $R_{\mathrm{lin}}$ & C4 \\
\hline
$N=14,R=5$ half-fill & $30$ & $0.99951$ & $0.99945$ & TIE \\
audit $N=16,R=5$ & $15$ & $0.889$ & $0.892$ & TIE \\
Paper~29 $E<0$, $R=4$ & $294$ & $0.808$ & $0.810$ & TIE \\
\hline
\end{tabular}
\end{center}

\begin{admission}
I changed the filling so C4 could be scored. It is a tie. I did
not call half-filling the Planck vacuum, and I did not write
that the parabola won.
\end{admission}

\end{document}
